\documentclass[10pt, twocolumn, landscape, a4paper]{article}

\usepackage[margin=1.57cm]{geometry}
\usepackage[utf8]{inputenc}
\usepackage[spanish,es-noshorthands]{babel}
\usepackage{amssymb, amsmath, amsbsy}
\usepackage{verbatim}
\usepackage{mathpazo}
\usepackage{tasks}
\usepackage{colortbl}
\usepackage{enumerate}
\usepackage{graphicx}
\usepackage{enumitem}
\usepackage[pdftex]{hyperref}
\hypersetup{pdftitle={Lógica},pdfauthor={Luis Carrillo Gutiérrez}}

\renewcommand{\familydefault}{\sfdefault}
\setlength{\columnsep}{1.2cm}
\pagestyle{empty} 

\begin{document}
\section*{Lógica Proposicional}

\subsection*{Lógica proposicional}

\noindent La lógica proposicional llamada también \emph{lógica simbólica} o \emph{lógica matemática}, es parte de la lógica que estudia las \textbf{proposiciones} y la relación que existe entre ellas a través de las \textbf{variables proposicionales} y los \textbf{conectores}.

\subsubsection*{Enunciado}

\noindent Se denomina enunciado a toda expresión del lenguaje común, el cual puede ser una frase, oración o expresión algebraica. \\

\noindent Ejemplos :
\begin{itemize}[noitemsep]
\item{¿Qué dia es hoy?}
\item{¡Auxilio!}
\item{Lima es la capital del Perú}
\item{$x + 2 > 3$}
\end{itemize}

\subsubsection*{Proposición lógica}
\noindent Llamada también \emph{enunciado cerrado}; es todo enunciado que tiene la cualidad de ser {\tt verdadero} o {\tt falso} sin ninguna ambigüedad. Estas se denotan con letras minúsculas tales como : $p$, $q$, $r$, $s$, $\ldots$ . \\
% Las proposiciones lógicas

\noindent Ejemplos :
\begin{itemize}[noitemsep]
\item{El sol es una estrella}
\item{2 + 3 > 6}
\item{Luis y María son hermanos}
\item{Juan estudia y trabaja} 
\end{itemize}

\paragraph*{Proposiciones simples} -- Llamadas también \emph{proposiciones atómicas} o elementales; son aquellas proposiciones con una sola idea, carecen de conjunciones gramaticales o del adverbio de negación (no). \\

\noindent Ejemplos :
\begin{itemize}[noitemsep]
\item{$p$ : Rubén es arquitecto}
\item{$q$ : Luis es compañero de José}
\item{$r$ : 3 y 4 son números consecutivos} 
\end{itemize}

\paragraph*{Proposiciones compuestas} -- Llamados también \emph{proposiciones moleculares}; son aquellas proposiciones con dos o más ideas unidas por conjunciones gramaticales; o que contienen el adverbio de negación (no).

\noindent Ejemplos :
\begin{itemize}[noitemsep]
\item{Luciana estudia contabilidad o administración}
\item{$2 + 8 \ge 6$}
\item{No es cierto que hoy sea lunes}
\item{Estudia, entonces ingresarás a la universidad}
\end{itemize}

\subsubsection*{Conectores lógicos}
\noindent Son símbolos que reemplazan a las conjunciones gramaticales y al adverbio de negación; estos conectores permiten relacionar dos o más proposiciones simples, entre los más importantes tenemos :
\subparagraph*{Negación} -- Cuya notación es $\neg$, se lee como : \emph{No es cierto que ...}. Algunas equivalencias de la negación son : No, no es el caso,
Por ejemplo :
\begin{eqnarray*}
p &:& \text{Luis viajó a Ica}. \\
\neg p &:& \text{No es cierto que Luis viajó a Ica}.
\end{eqnarray*}

Tabla de verdad : \\

\begin{center}
\begin{tabular}{|c|c|}
	\arrayrulecolor{cyan}\hline
	$p$ & $\neg p$ \\
	\hline
	V & F \\
	\hline
	F & V \\
	\hline
\end{tabular}
\end{center}

\subparagraph*{Conjunción} -- Cuya notación es $\wedge$, se lee como : ... y ... Por ejemplo :
\begin{eqnarray*}
p &:& \text{Ángel estudia}. \\
q &:& \text{Ángel trabaja}. \\
p \wedge q &:& \underbrace{\text{Ángel estudia}}_{p}\quad \underbrace{y}_{\wedge}\quad \underbrace{trabaja}_{q}.
\end{eqnarray*}

Tabla de verdad (matriz principal): \\

\begin{center}
\begin{tabular}{|c|c|c|}
	\arrayrulecolor{cyan}\hline
	$p$ & $q$ &  $p \wedge q$ \\
	\hline
	V & V & V \\
	\hline
	V & F & F \\
	\hline
	F & V & F \\
	\hline
	F & F & F \\
	\hline
\end{tabular}
\end{center}

\end{document}
